\subsection{Cohomology of Algebras}

In this subsection, K is a commutative ring, A a K-algebra, B the K-algebra \(A \otimes_{K} A^{\circ}\) , and \(\varepsilon\) the K-linear mapping from B to A defined by \(\varepsilon(x \otimes y) = xy\) for x, y in A. For \(n \in N\) , we denote the tensor product over K of \(n + 2\) copies of the K-module A by \(B_{n}\) . We view it as an (A, A)-bimodule (and also as a B-module); we endow it with the structure of a left A-module deduced from the left A-module structure of the first factor of the tensor product and with the structure of a right A-module deduced from the right A-module structure of the last factor. In particular, \(B_{0}\) is just the (A, A)-bimodule B.

For every integer \(n \geqslant 1\) , we define homomorphisms of bimodules \(d_{n}^{i}\) for \(i \in [0, n]\) and \(d_{n}\) from \(B_{n}\) to \(B_{n-1}\) by the formulas

\[
d _ {n} ^ {i} (x _ {0} \otimes \dots \otimes x _ {n + 1}) = x _ {0} \otimes \dots \otimes x _ {i} x _ {i + 1} \otimes \dots \otimes x _ {n + 1}\tag{4}
\]

for \(i\in [0,n]\) and

\[
d _ {n} = \sum_ {i = 0} ^ {n} (- 1) ^ {i} d _ {n} ^ {i}.\tag{5}
\]

We denote the mapping \(\varepsilon : \mathrm{B}_0 \to \mathrm{A}\) by \(d_0 = d_0^0\) .

Let \(n\) be an integer \(\geqslant 1\) . For \(0 \leqslant i < j \leqslant n\) , we have

\[
d _ {n - 1} ^ {i} \circ d _ {n} ^ {j} = d _ {n - 1} ^ {j - 1} \circ d _ {n} ^ {i},\tag{6}
\]

and from this we deduce

\[
\begin{array}{c} d _ {n - 1} \circ d _ {n} = \sum_ {0 \leqslant i <   j \leqslant n} (- 1) ^ {i + j} d _ {n - 1} ^ {i} \circ d _ {n} ^ {j} + \sum_ {0 \leqslant j \leqslant i \leqslant n - 1} (- 1) ^ {i + j} d _ {n - 1} ^ {i} \circ d _ {n} ^ {j} \\ = \sum_ {0 \leqslant i <   j \leqslant n} (- 1) ^ {i + j} d _ {n - 1} ^ {j - 1} \circ d _ {n} ^ {i} + \sum_ {0 \leqslant j \leqslant i \leqslant n - 1} (- 1) ^ {i + j} d _ {n - 1} ^ {i} \circ d _ {n} ^ {j} \end{array}
\]

and consequently

\[
d _ {n - 1} \circ d _ {n} = 0.\tag{7}
\]

Let P be an (A,A)-bimodule. For any integer \(n \geqslant 0\) , we denote the K-module of K-multilinear mappings from \(A^{n}\) to P by \(C^{n}(A,P)\) . The mapping \(\alpha^{n}: C^{n}(A,P) \to \operatorname{Hom}_{\mathrm{B}}(\mathrm{B}_{n},\mathrm{P})\) that sends \(f \in \mathrm{C}^{n}(A,P)\) to the homomorphism \(\alpha^{n}(f)\) characterized by

\[
\alpha^ {n} (f) \left(x _ {0} \otimes \dots \otimes x _ {n + 1}\right) = x _ {0} f (x _ {1}, \ldots , x _ {n}) x _ {n + 1}\tag{8}
\]

is an isomorphism of K-modules.

We denote by \(\partial^{n}\) (for \(n \geqslant 0\) ) the unique K-linear mapping from \(\mathrm{C}^{n}(\mathrm{A},\mathrm{P})\) to \(\mathrm{C}^{n+1}(\mathrm{A},\mathrm{P})\) that makes the following diagram commutative:

\[
\begin{array}{c} \mathrm{C} ^ {n} (\mathrm{A}, \mathrm{P}) \xrightarrow {\partial^ {n}} \mathrm{C} ^ {n + 1} (\mathrm{A}, \mathrm{P}) \\ \Biggl \downarrow_ {\alpha^ {n}} \\ \mathrm{Hom} _ {\mathrm{B}} (\mathrm{B} _ {n}, \mathrm{P}) \xrightarrow {\mathrm{Hom} (d _ {n + 1} , 1 _ {\mathrm{P}})} \mathrm{Hom} _ {\mathrm{B}} (\mathrm{B} _ {n + 1}, \mathrm{P}). \end{array}
\]

By definition, we therefore have

\[
(\alpha^ {n + 1} \circ \partial^ {n}) (f) = \alpha^ {n} (f) \circ d _ {n + 1}\tag{9}
\]

for every \(f \in \mathrm{C}^n(\mathrm{A},\mathrm{P})\) . In other words, we have

\[
\partial^ {n} (f) \left(x _ {0}, \dots , x _ {n}\right) = \alpha^ {n} (f) \left(d _ {n + 1} \left(1 \otimes x _ {0} \otimes \dots \otimes x _ {n} \otimes 1\right)\right)
\]

for \(x_0, \ldots, x_n\) in A and \(f\) in \(\mathrm{C}^n(\mathrm{A}, \mathrm{P})\) , that is,

\[
\begin{array}{r l} & {\partial^ {n} (f) (x _ {0}, \ldots , x _ {n}) = x _ {0} f (x _ {1}, \ldots , x _ {n})} \\ & {\qquad + \sum_ {i = 0} ^ {n - 1} (- 1) ^ {i + 1} f (x _ {0}, \ldots , x _ {i - 1}, x _ {i} x _ {i + 1}, x _ {i + 2}, \ldots , x _ {n})} \\ & {\qquad \qquad \qquad \qquad \qquad + (- 1) ^ {n + 1} f (x _ {0}, \ldots , x _ {n - 1}) x _ {n}.} \end{array}\tag{10}
\]

By (7) and (9), we have

\[
\partial^ {n + 1} \circ \partial^ {n} = 0 \quad \text {   for   every   } n \geqslant 0.\tag{11}
\]

We denote the K-module \(\operatorname{Ker} \partial^0\) by \(\mathrm{H}^0(\mathrm{A},\mathrm{P})\) and, for \(n \geqslant 1\) , the K-module \(\operatorname{Ker} \partial^n/\operatorname{Im} \partial^{n-1}\) by \(\mathrm{H}^n(\mathrm{A},\mathrm{P})\) . We identify the K-module \(\mathrm{C}^0(\mathrm{A},\mathrm{P})\) with \(\mathrm{P}\) , and we have \(\mathrm{C}^1(\mathrm{A},\mathrm{P}) = \operatorname{Hom}_\mathrm{K}(\mathrm{A},\mathrm{P})\) . The mappings \(\partial^n\) for \(n \leqslant 2\) are given by the formulas

\[
\partial^ {0} (p) (a) = a p - p a \qquad \mathrm{for\ every} p \in \mathrm{P},\tag{12}
\]

\[
\partial^ {1} (f) \left(a, a ^ {\prime}\right) = a f \left(a ^ {\prime}\right) - f \left(a a ^ {\prime}\right) + f (a) a ^ {\prime} \quad \text { for } f \in \mathrm{C} ^ {1} (\mathrm{A}, \mathrm{P}),\tag{13}
\]

\[
\partial^ {2} (f) \left(a, a ^ {\prime}, a ^ {\prime \prime}\right) = a f \left(a ^ {\prime}, a ^ {\prime \prime}\right) - f \left(a a ^ {\prime}, a ^ {\prime \prime}\right) + f \left(a, a ^ {\prime} a ^ {\prime \prime}\right) - f \left(a, a ^ {\prime}\right) a ^ {\prime \prime}\tag{14}
\]

for \(f \in \mathrm{C}^2(\mathrm{A},\mathrm{P})\) .

So \(H^{0}(A,P)\) is the K-submodule of P consisting of the elements p such that ap = pa for every \(a \in A\) , and \(H^{1}(A,P)\) is the quotient of the K-module \(\mathrm{Der}_{\mathrm{K}}(\mathrm{A},\mathrm{P})\) of K-derivations from A to P (III, §10, No.~2, p.~553) by the K-submodule consisting of the derivations of the form \(a \mapsto ap - pa\) with \(p \in P\) (called inner derivations).

